% =============================================================================
% 02-pengesahan.tex — Halaman Pengesahan
% =============================================================================
\clearpage
\thispagestyle{empty}

\begin{singlespace}
\begin{center}
  {\fontsize{14}{16}\selectfont\bfseries HALAMAN PENGESAHAN}\\[0.5cm]
  {\fontsize{14}{16}\selectfont\bfseries LAPORAN KERJA PRAKTIK}\\[1.5cm]
  {\fontsize{12}{14}\selectfont\bfseries \judulLaporan}\\[1cm]
\end{center}

\noindent Laporan Kerja Praktik ini telah dipertahankan di hadapan tim penguji pada:

\vspace{0.5cm}
\noindent
\begin{tabular}{@{}ll@{}}
  Hari/Tanggal & : \rule{6cm}{0.4pt} \\[0.3cm]
  Tempat & : Program Studi Informatika, FTC, UKRIDA \\
\end{tabular}

\vspace{1cm}

\noindent dan dinyatakan \textbf{LULUS} dengan nilai \rule{3cm}{0.4pt}.

\vspace{1.5cm}

\begin{center}
  {\bfseries Tim Penguji}
\end{center}

\vspace{0.5cm}

\noindent
\begin{minipage}[t]{0.48\textwidth}
  \begin{center}
    Penguji I,\\[2.5cm]
    \rule{5cm}{0.4pt}\\
    NIP.
  \end{center}
\end{minipage}
\hfill
\begin{minipage}[t]{0.48\textwidth}
  \begin{center}
    Penguji II,\\[2.5cm]
    \rule{5cm}{0.4pt}\\
    NIP.
  \end{center}
\end{minipage}

\vspace{1cm}

\begin{center}
  Mengetahui,\\
  Ketua Program Studi Informatika\\[2.5cm]
  \rule{6cm}{0.4pt}\\
  NIP.
\end{center}

\end{singlespace}
